\documentclass[11pt,a4paper]{article}

\usepackage[T1]{fontenc}
\usepackage[margin=2.5cm]{geometry}
\usepackage{amsmath}
\usepackage{array}
\usepackage{hyperref}

\title{Native LM Backend: Project Change Record}
\date{3 September 2026}

\begin{document}

\maketitle

\section{Purpose of This Record}

This document records every project-specific change made while integrating the
native Languasco--Migliardi (LM) log-likelihood implementation into LM-IDNet.
It is intended to make the integration reviewable and to provide a concise
account that can be reported to the authors of the native implementation.

The original reference copy under \path{lm-algorithm/code/} is kept unchanged.
All integration edits are made to the runtime project copy under
\path{src/lm_idnet/algorithms/native/}.

\section{Baseline Before Native Changes}

Before modifying the native source, the following checks were completed:

\begin{itemize}

\item The focused likelihood and estimator-factory tests passed: 5 passed.

\item The normal project suite passed: 195 passed and 10 tests were deselected
by the configured test markers.

\item The runtime C sources, headers, and coefficient tables were compared with
the reference capsule.

\item The two copies of \path{LM_time_lib.c} were identical before the first
edit, with SHA-256 digest:

\begin{verbatim}
306b33aae05db49e519909c36e91d25112b5fce75a3fb144858ed0d727f22d66
\end{verbatim}

\end{itemize}

\section{Change 001: Exact Handling of a Zero Count}

\subsection{Status}

Implemented and directly validated against a temporary shared library compiled
from the modified source.

\subsection{File Changed}

\begin{verbatim}
src/lm_idnet/algorithms/native/LM_time_lib.c
\end{verbatim}

The corresponding reference file was not changed:

\begin{verbatim}
lm-algorithm/code/LM_time_lib.c
\end{verbatim}

After this change, the digests are:

\begin{verbatim}
Project copy:
e948740488ed6ef731a7eabc7f02dcdbd83f666fcfbd793010ac14003bdf5e15

Reference copy:
306b33aae05db49e519909c36e91d25112b5fce75a3fb144858ed0d727f22d66
\end{verbatim}

\subsection{Problem Being Addressed}

The function \texttt{LM\_logL} evaluates a paired log-gamma difference. In the
application parameterization, its mathematical contribution is of the form

\[
D(\alpha,y)=\log\Gamma(\alpha+y)-\log\Gamma(\alpha),
\]

where \(y\) is a non-negative packet count. Network-traffic windows commonly
contain categories for which \(y=0\). A completely silent window also has a
total count \(N=0\).

For zero counts, the exact value is

\[
D(\alpha,0)
=\log\Gamma(\alpha)-\log\Gamma(\alpha)
=0.
\]

The supplied implementation already contained exact special cases for
\(y=1\) and \(y=2\), but it did not contain the corresponding base case for
\(y=0\). Sending zero through the general Euler--Maclaurin expression is
unnecessary and does not directly express this exact identity.

\subsection{Applied Modification}

The following guard was inserted at the beginning of the existing special-value
section in \texttt{LM\_logL}, before the cases for one and two:

\begin{verbatim}
if (y == 0)
    {
    return 0.0;
    }
\end{verbatim}

The semantic source diff is:

\begin{verbatim}
 /* special values */
+if (y == 0)
+    {
+    return 0.0;
+    }
 if (y == 1)
\end{verbatim}

The patching tool also normalized one redundant blank line at the end of the C
file. That whitespace-only difference has no effect on compilation or behavior.

\subsection{Why the Change Was Made in This Function}

Handling the case in \texttt{LM\_logL} covers both places where a zero can
occur:

\begin{itemize}

\item a category contribution with \(x_k=0\); and

\item the concentration contribution for a completely silent row with
\(N=0\).

\end{itemize}

Placing the guard in the shared paired-difference function avoids duplicating
zero checks in the future matrix wrapper or in Python. It also ensures that any
other caller of \texttt{LM\_logL} receives the same exact behavior.

\subsection{Effect on the LM Numerical Method}

This change does not modify:

\begin{itemize}

\item the Euler--Maclaurin expansion;

\item coefficient values;

\item truncation-order selection;

\item error-bound calculations;

\item horizontal shifting; or

\item results for positive counts.

\end{itemize}

It only returns a known exact identity before the approximation is entered.

\subsection{Native Verification}

The host development environment did not contain a C compiler. The modified
source was therefore compiled using the existing \texttt{lm-idnet:latest}
Docker image, which is the supported project toolchain. The shared object was
written only to a temporary directory; the versioned project binary was not
replaced.

The compiler invocation was equivalent to:

\begin{verbatim}
gcc -O3 -fPIC -shared LM_time_lib.c -lm -o LM_time_lib.so
\end{verbatim}

The compiled \texttt{LM\_logL} function was then called through
\texttt{ctypes} with three different positive reciprocal parameters, expansion
orders, and horizontal-shift values. All three calls returned exactly zero:

\begin{verbatim}
zero-count values: [0.0, 0.0, 0.0]
\end{verbatim}

A four-category silent row was also evaluated by combining the concentration
term and four category terms in the same form as the native likelihood:

\begin{verbatim}
silent-row kernel: 0.0
\end{verbatim}

The existing focused Python tests were rerun after the source edit:

\begin{verbatim}
tests/test_likelihood.py          4 passed
tests/test_estimator_factory.py   1 passed
Total                             5 passed
\end{verbatim}

The normal project suite was also rerun after the native change:

\begin{verbatim}
195 passed, 10 deselected
\end{verbatim}

Finally, \texttt{git diff --check} reported no whitespace errors.

\subsection{Current Limitation}

This change has been tested directly at the native function boundary, but the
public Python LM backend is still a placeholder. An end-to-end Python test for
silent matrices will be added when the native adapter and matrix entry point are
implemented. Until then, the modified shared object is compiled only for direct
verification and is not used by normal application commands.

\subsection{Short Report for the Native-Code Authors}

The project copy adds an exact \(y=0\) base case to
\texttt{LM\_logL}. The mathematical paired log-gamma difference is exactly zero
for this input. The guard is necessary for sparse protocol vectors and fully
silent IoT traffic windows. It does not alter any positive-count calculation or
any part of the LM approximation and error-control method. The original capsule
source remains unchanged.

\section{Change 002: Explicit Paths and Checked Coefficient Loading}

\subsection{Status}

Implemented and validated at the native shared-library boundary. The Python adapter that
will translate the returned statuses into application exceptions remains a later
integration step.

\subsection{Files Changed}

\begin{verbatim}
src/lm_idnet/algorithms/native/LM_time_lib.c
src/lm_idnet/algorithms/native/LM_hor_shift.h
\end{verbatim}

The corresponding files under \path{lm-algorithm/code/} were not changed.

\subsection{Problem Being Addressed}

The supplied \texttt{initBern\_logL} function opened the two coefficient files by
bare filename. Its success therefore depended on the process current working
directory. It also wrote through allocated pointers without checking whether
the files opened, whether allocation succeeded, or whether all coefficients
were read. Used as a library, these failure modes could cause invalid state or
a process crash.

\subsection{Applied Modification}

A checked initializer was added and declared in \path{LM_hor_shift.h}:

\begin{verbatim}
int initBern_logL_paths(
    const char *bernoulli_path,
    const char *error_path
);
\end{verbatim}

This function:

\begin{enumerate}

\item rejects null path pointers;

\item opens the two caller-supplied paths;

\item allocates temporary 100-element arrays;

\item verifies all 100 reads from each file;

\item closes both files on every path after they have been opened;

\item preserves the previously loaded tables if opening, allocation, or
parsing fails; and

\item replaces and releases the old tables only after both new tables have
been read successfully.

\end{enumerate}

The return statuses are:

\begin{center}
\begin{tabular}{>{\ttfamily}l p{10cm}}
0 & both tables loaded successfully; \\
1 & a path was null or a file could not be opened; \\
2 & memory allocation failed; \\
3 & a file did not contain all 100 parseable coefficients. \\
\end{tabular}
\end{center}

The original no-argument \texttt{initBern\_logL} symbol remains available for
compatibility. It delegates to the checked function with the original relative
filenames and writes one diagnostic to standard error if loading fails. The
new Python integration will call \texttt{initBern\_logL\_paths} with absolute
paths instead.

Only the following tables are loaded by this likelihood-specific initializer:

\begin{verbatim}
bernreal-norm-100.txt
err_coeff-100.txt
\end{verbatim}

No coefficient value or file was modified.

\subsection{Native Verification}

The modified library was loaded after changing the test process current
directory to \path{/tmp}. Initialization with absolute paths succeeded. The
following failure cases returned the expected nonzero status without crashing:

\begin{itemize}

\item two null paths: status 1;

\item a missing file: status 1; and

\item a nonnumeric file in place of a coefficient table: status 3.

\end{itemize}

After a valid load, an attempted malformed reload returned status 3, and a
subsequent likelihood calculation still succeeded. This confirms that a failed
reload does not destroy the previously valid coefficient state.

The initializer was also called successfully 1,000 times. Process resident
memory remained at 11,644 KiB before and after those reloads in this stress
check.

\subsection{Effect on the LM Numerical Method}

This change affects file location, validation, and table ownership only. It
does not change the coefficient files, their in-memory values, or any LM
formula.

\subsection{Short Report for the Native-Code Authors}

The project copy adds an explicit-path, status-returning initializer for the two
tables used by the log-likelihood path. It reads into temporary storage and
commits the replacement only after both complete tables have been validated.
The original initializer remains as a compatibility wrapper. The change removes
the current-working-directory dependency and makes initialization failures
recoverable by a library caller.

\section{Change 003: Explicit Workspace Ownership and Allocation Checks}

\subsection{Status}

Implemented for every allocation exercised by the current parameter and
log-likelihood path. Final cleanup after a matrix calculation will be connected
when the matrix-level entry point is added; that entry point does not yet exist.

\subsection{Problem Being Addressed}

The native implementation retained four parameter arrays in global pointers.
Every call to \texttt{init\_params} replaced those pointers with new allocations
without releasing the preceding arrays. In addition,
\texttt{opterrNoPrint\_logL} allocated two temporary arrays on every likelihood
calculation and did not release them. Repeated fitting iterations would
therefore continually increase native memory use.

The allocation results were not checked before their pointers were used.

\subsection{Applied Modifications}

A private \texttt{free\_params} helper now releases the arrays referenced by
\texttt{X}, \texttt{probs}, \texttt{psioverprobs}, and \texttt{logprobs}, then
sets all four pointers to \texttt{NULL}. The parameter initializer invokes it
before allocating a replacement workspace and again if any replacement
allocation fails.

\texttt{init\_params} now rejects a null count pointer, a nonpositive category
count, or a nonpositive precision value by returning \texttt{-1}. It also
returns \texttt{-1} if any of its four allocations fails. No pointer is written
through after a failed allocation.

Within \texttt{opterrNoPrint\_logL}, the two category-length temporary arrays
were replaced by one allocation of \texttt{2*K} doubles, divided into two
non-overlapping halves. This keeps the same values and indexing while requiring
only one allocation check. The block is freed on both existing return paths:

\begin{itemize}

\item failure to meet the requested accuracy within the horizontal-shift
limit; and

\item successful selection of the expansion order and horizontal shift.

\end{itemize}

If that temporary allocation fails, or if the requested accuracy cannot be met
within the existing horizontal-shift limit, \texttt{opterrNoPrint\_logL}
returns \texttt{NULL}. \texttt{loggamma\_LM} checks that return and propagates
the failure as \texttt{NAN}. The future matrix wrapper can therefore reject the
non-finite value rather than mistaking the failure for a valid likelihood
result.

\texttt{loggamma\_LM} also invokes \texttt{free\_params} before returning
its existing precision-limit sentinel. Thus, both failure returns currently
visible in that function release the active parameter workspace.

The coefficient arrays have a separate, longer lifetime. The private
\texttt{freeBern\_logL} helper releases and nulls those two pointers only when
a complete replacement pair has been loaded successfully.

\subsection{Ownership After This Change}

\begin{center}
\begin{tabular}{p{4.2cm} p{5.3cm} p{4.0cm}}
\textbf{Storage} & \textbf{Lifetime} & \textbf{Released by} \\
\hline
Coefficient tables & Across likelihood calls, until a successful reload &
\texttt{freeBern\_logL} before replacement \\
Parameter arrays & From successful \texttt{init\_params} until replacement or
the future matrix boundary & \texttt{free\_params} \\
Error-search workspace & One \texttt{opterrNoPrint\_logL} call & The function
itself on success or failure \\
\end{tabular}
\end{center}

\subsection{Deliberately Deferred Connection Point}

There is not yet a matrix-level native function, so there is currently no
``last row'' boundary at which to invoke \texttt{free\_params}. The helper is
private in the same C translation unit and is ready to be called from the
matrix wrapper in the next implementation step. Until then, the most recent
parameter workspace remains live for the next \texttt{loggamma\_LM} call and is
released before the next \texttt{init\_params} call. This preserves the calling
sequence expected by the supplied API without introducing a speculative public
cleanup API.

\subsection{Native Verification}

The following checks were made through \texttt{ctypes} against a temporary
shared library:

\begin{itemize}

\item valid initialization and likelihood evaluation produced the finite
value \texttt{-8.923657799857487};

\item null counts, zero categories, and zero precision were rejected with
\texttt{-1};

\item the existing precision-limit path returned its documented value of 99,
after which parameter initialization and a normal likelihood call succeeded;

\item 20,000 alternating three-category and 256-category initializations
increased resident memory by only 68 KiB (11,576 to 11,644 KiB); and

\item 10,000 repeated 256-category likelihood evaluations returned the same
finite value each time while resident memory remained at 11,644 KiB.

\end{itemize}

The small resident-memory differences are consistent with allocator and process
bookkeeping; importantly, growth was not proportional to the number of calls.

\subsection{Effect on the LM Numerical Method}

The contents, length, and indexing of the working arrays are unchanged. No
formula, loop termination condition, coefficient, error bound, expansion order,
or horizontal-shift rule was changed.

\subsection{Short Report for the Native-Code Authors}

The project copy now gives the parameter arrays and per-likelihood error-search
arrays explicit owners. Existing arrays are released before replacement,
temporary arrays are released on all returns, pointers are nulled after release,
and allocations are checked. Repeated calls no longer leak these workspaces.
The final parameter cleanup call will sit in the forthcoming matrix wrapper,
where the end of a calculation is unambiguous.

\section{Consolidated Verification After Changes 001--003}

The source was compiled using the existing project Docker image with
\texttt{gcc -O3 -Wall -Wextra -fPIC -shared} and linked with
\texttt{-lm}. Compilation succeeded. The reported warnings concern unused
parameters or variables already present in the supplied implementation; no
invalid declaration, missing symbol, or linkage warning was reported.

The Python regression results after Changes 001--003 are:

\begin{verbatim}
Focused likelihood/factory tests: 5 passed
Normal project suite:             195 passed, 10 deselected
git diff --check:                 passed
\end{verbatim}

The final SHA-256 digests at this checkpoint are:

\begin{verbatim}
Runtime LM_time_lib.c:
8ddea61941d64977562c640c8264d75d5127ff353150240be4c019f9fe6b6c3d

Reference LM_time_lib.c (unchanged):
306b33aae05db49e519909c36e91d25112b5fce75a3fb144858ed0d727f22d66

Runtime LM_hor_shift.h:
1ce8a17c4a50057f4d92a6cbf580045320ba221d6b4883236f8bbbbb2d6580fa

Reference LM_hor_shift.h (unchanged):
3ba1ac510b0fba74a64ae4bc63be5301c214ee783a2fbde0e883188f32a3be96

bernreal-norm-100.txt (unchanged):
5072a78dbf4af4e317704e714d25fa4ea559bbb169d2652a60328381553e09dd

err_coeff-100.txt (unchanged):
48eb77cf4c9527a204022899469a1877fd667fc53d38ab164fabf1a0d6adebf6
\end{verbatim}

The versioned \path{LM_time_lib.so} was not replaced. All native checks used
a temporary shared library, so the normal application still uses its existing
backend behavior until the remaining integration steps are completed.

\end{document}
